\clearpage
\section{Анализ роли модульного тестирования}

\subsection{Задачи, решаемые модульными тестами}

Модульные тесты решают задачу автоматизированной проверки небольших, логически завершённых частей программной системы. В выполненных работах единицей проверки являлся отдельный метод или небольшая группа взаимосвязанных свойств класса. Такой масштаб позволяет установить непосредственную связь между входными данными, выполненной операцией и наблюдаемым результатом.

Первая функция тестов состоит в проверке соответствия реализации требованиям. Например, тесты метода \texttt{IsPalindrome} формализуют правила игнорирования регистра, пробелов и пунктуации. Без тестов эти положения остаются только текстовым описанием; с тестами они превращаются в исполняемую спецификацию. Если реализация перестаёт удовлетворять хотя бы одному правилу, система тестирования указывает конкретный нарушенный сценарий.

Вторая функция состоит в предотвращении регрессий. Изменение внутреннего алгоритма, обновление зависимостей или добавление нового задания не должно нарушать ранее реализованное поведение. Запуск всех 32 тестов после каждой сборки обеспечивает повторную проверку заданий 01--05. Следовательно, тесты создают контролируемые условия для рефакторинга: разработчик может менять структуру кода, сохраняя внешний контракт.

Третья функция связана с локализацией дефектов. Если единый большой сценарий завершается ошибкой, поиск причины может потребовать анализа множества компонентов. Небольшой тест, посвящённый одному свойству, существенно сужает область поиска. Так, отказ теста \codebreak{GetReverseEnumerator_ReturnsItemsInReverseOrder} непосредственно указывает на ошибку обратного перечисления, а не на методы добавления или фильтрации.

Четвёртая функция заключается в документировании поведения на примерах. Название теста, его подготовительная часть и утверждение показывают способ использования программного интерфейса. Это особенно заметно в заданиях с итераторами и рефлексией: тесты демонстрируют момент материализации последовательности, состав ожидаемых метаданных и реакцию на отсутствие метода.

Наконец, модульные тесты являются основой автоматической проверки в CI. Локально успешный результат сам по себе не гарантирует воспроизводимость в другой среде. GitHub Actions восстанавливает зависимости, собирает решение и запускает тесты на отдельном исполнителе. Успешное прохождение этой процедуры подтверждает, что результат не зависит от случайного состояния рабочего компьютера.

\subsection{Влияние разработки через тесты на скорость работы}

Разработка через тесты может как замедлять отдельный этап, так и ускорять полный цикл создания и сопровождения программы. На начальном этапе требуется сформулировать сценарии, подготовить данные и определить ожидаемые результаты. Для небольшого учебного метода написание теста иногда занимает сопоставимое с реализацией время. Поэтому при оценке только длительности первичного набора кода тестирование воспринимается как дополнительная работа.

Однако такая оценка не учитывает время последующей отладки. Формулирование теста вынуждает заранее уточнить контракт метода. В задании о палиндроме это позволило определить поведение пустых и очищаемых строк до появления скрытых ошибок. В задании о рефлексии тест параметров зафиксировал, что результат должен включать не только имена параметров, но и возвращаемый тип. Чем раньше обнаружена неоднозначность, тем дешевле её устранение.

Тесты ускоряют повторные изменения. После добавления нового проекта общее решение собиралось и проверялось полностью. Ручная проверка всех пяти заданий потребовала бы многократного ввода данных и визуального сравнения результатов. Автоматический запуск выполняет те же проверки воспроизводимо и за ограниченное время. Экономия возрастает с числом изменений и продолжительностью жизненного цикла проекта.

Поэтому обоснованным является вывод, что разработка через тесты умеренно замедляет первое написание простейшей версии, но ускоряет получение устойчивого результата. Для одноразового прототипа выигрыш может быть невелик. Для кода, который развивается, проверяется преподавателем, переносится между средами и подвергается рефакторингу, суммарные затраты времени снижаются благодаря раннему выявлению дефектов и автоматизации регрессионной проверки.

\subsection{Основания идеи модульных тестов}

Идея модульного тестирования основана на декомпозиции сложной системы. Программа рассматривается как совокупность компонентов с определёнными контрактами. Если каждый компонент получает контролируемые входные данные и формирует ожидаемый результат, то вероятность корректного взаимодействия компонентов повышается. Модульный тест не доказывает отсутствие всех возможных ошибок, но проверяет существенные свойства на выбранном множестве случаев.

Вторым основанием является принцип наблюдаемости. Поведение программы должно быть выражено через результат, изменение состояния или зафиксированное взаимодействие. Именно поэтому в задании об интерфейсах методы движения и выстрела связаны с инкапсулированным состоянием. Проверка только факта отсутствия исключения была бы слабой: она не подтверждала бы, что команда действительно оказала требуемый эффект.

Третьим основанием выступает воспроизводимость эксперимента. Каждый тест задаёт исходное состояние, выполняет одно действие и сравнивает наблюдение с ожидаемым постусловием. При одинаковой версии кода результат должен быть одинаковым независимо от числа запусков. Такая организация сближает тестирование с экспериментальным методом: гипотеза о корректности конкретного свойства подвергается повторяемой проверке.

Четвёртым основанием является быстрая обратная связь. Чем меньше временной интервал между внесением изменения и обнаружением нарушения, тем проще связать дефект с его причиной. Автоматический запуск тестов при отправке изменений в репозиторий создаёт единый контрольный контур: изменение кода, сборка, проверка и получение результата CI.

Следовательно, модульное тестирование основано на сочетании декомпозиции, явного контракта, наблюдаемости, изоляции и воспроизводимой проверки. Его практическая ценность определяется не количеством утверждений само по себе, а тем, насколько выбранные сценарии отражают значимые свойства программы и возможные источники ошибок.

\subsection{Ограничения выполненного тестирования}

Несмотря на успешное выполнение всех тестов, покрытие не является исчерпывающим доказательством корректности. Набор ориентирован на требования заданий и наиболее важные граничные случаи. Например, метод работы с палиндромами можно дополнительно проверить на строках с цифрами и символами разных алфавитов; сервис студентов --- на пустой коллекции и отсутствии оценок; генератор --- на нулевом и отрицательном значении \texttt{count}; анализатор --- на перегруженных методах.

Указанные случаи не отменяют ценность существующих проверок, но показывают необходимость риск-ориентированного расширения набора тестов. При дальнейшем развитии проекта новые тесты должны добавляться для каждого обнаруженного дефекта и каждого уточнённого требования. Такой подход превращает набор тестов в накапливаемую базу знаний о поведении системы.
